% ===== 原稿 PDF 第 30 页 =====
% 矩阵乘法的维数关系、第 (2)(3) 条、例、乘法的性质及结合律的证明、方阵的乘方（定义 7）
% 本页是「矩阵的运算」一节的续页，原稿无灰底方框节标题块，故不新开 \section

$AB$ 的行数 $=$ $A$ 的行数，$AB$ 的列数 $=$ $B$ 的列数。

\begin{center}
\begin{tikzpicture}[>=Stealth,every node/.style={inner sep=2pt}]
  \node (A)   at (0,0)      {$A$};
  \node (B)   at (3.0,0)    {$B$};
  \node (AB)  at (6.4,0)    {$AB$};
  \node (Ad)  at (0,-0.6)   {$m\times n$};
  \node (Bd)  at (3.0,-0.6) {$n\times p$};
  \node (ABd) at (6.4,-0.6) {$m\times p$};
  \draw[->] ([xshift=6mm]B.east) -- ([xshift=-6mm]AB.west);
  \draw ([xshift=1mm,yshift=-1mm]Ad.south east) to[bend right=30] ([xshift=-1mm,yshift=-1mm]Bd.south west);
\end{tikzpicture}
\end{center}
% 原稿图中 A 下方标 m×n（有涂改），B 下方标 n×p（首字笔迹涂改不清），AB 下方标 m×p 处笔迹作 m×n；
% 此处按可乘条件及下方例子（2×3、3×4、2×4）与后文 B=(b_{ij})_{n×p} 取 A:m×n、B:n×p、AB:m×p
% 原稿在 A、B 的阶数之间画弧线相连（表示内阶数 n 相同），已用下弧示意

(2) $AB$ 的第 $(i,j)$ 元素 $=$ $A$ 的第 $i$ 行与 $B$ 的第 $j$ 列对应元素乘积之和。
% 原稿作「对应元素之和」（脱「乘积」二字）

\begin{example}
% 原稿在三个矩阵上方分别标注 2×3、3×4、2×4
\[
\begin{pmatrix}
1 & 0 & 1\\
2 & 1 & 0
\end{pmatrix}
\begin{pmatrix}
1 & 0 & 1 & 1\\
1 & 1 & 2 & -1\\
-1 & 0 & -1 & 0
\end{pmatrix}
=
\begin{pmatrix}
0 & 0 & 0 & 1\\
3 & 1 & 4 & 1
\end{pmatrix}
\]
\end{example}

(3) 矩阵乘法一般不满足交换律。

即使 $AB$，$BA$ 都有意义，但 $AB\ne BA$。

\begin{enumerate}[label=(\arabic*)]
\item $A=(1\ 1\ 1)$，$B=\begin{pmatrix}1\\1\\1\end{pmatrix}$，$AB=(3)$，
\[
BA=\begin{pmatrix}
1 & 1 & 1\\
1 & 1 & 1\\
1 & 1 & 1
\end{pmatrix}\ne AB;
\]
\item $A=\begin{pmatrix}0 & 1\\0 & 0\end{pmatrix}$，$B=\begin{pmatrix}0 & 0\\0 & 1\end{pmatrix}$，
\[
AB=\begin{pmatrix}0 & 1\\0 & 0\end{pmatrix}\ne BA=\begin{pmatrix}0 & 0\\0 & 0\end{pmatrix}.
\]
\end{enumerate}

\textbf{乘法的性质}
\begin{enumerate}[label=(\arabic*)]
\item 结合律：$(A\cdot B)\cdot C=A\cdot(B\cdot C)$；
\item 分配律：$(A+B)\cdot C=AC+BC$，
\[
A\cdot(B+C)=AB+AC;
\]
\item 与数乘的相容性：$c\cdot(A\cdot B)=(c\cdot A)\cdot B=A\cdot(c\cdot B)$；
\item 单位元：记 $A_{m\times n}$，$I_m\cdot A=A=A\cdot I_n$。
\end{enumerate}

\begin{proof}
只证 (1)。设 $A=(a_{ij})_{m\times n}$，$B=(b_{ij})_{n\times p}$，$C=(c_{ij})_{p\times q}$。

先考虑 $AB$ 的第 $(i,r)$ 元素：
\[
\sum_{k=1}^{n}a_{ik}b_{kr}.
\]
对 $(AB)\cdot C$ 的第 $(i,j)$ 元素：
% 原稿此行行首有一字（形似「对」），笔迹难辨，按上下文作「对」
\[
\sum_{r=1}^{p}\Big(\sum_{k=1}^{n}a_{ik}b_{kr}\Big)\cdot c_{rj}
=\sum_{r=1}^{p}\sum_{k=1}^{n}a_{ik}b_{kr}c_{rj}.
\]
$A(BC)$：先考虑 $BC$ 的第 $(k,j)$ 元素
% 原稿自「A(BC)」起用一个大左括号（括住下面三行），此处按行顺次排出
\[
\sum_{r=1}^{p}b_{kr}c_{rj},
\]
第 $(i,j)$ 元素
\[
=\sum_{k=1}^{n}a_{ik}\Big(\sum_{r=1}^{p}b_{kr}c_{rj}\Big)
=\sum_{k=1}^{n}\sum_{r=1}^{p}a_{ik}b_{kr}c_{rj}.
\]
\[
\Rightarrow\quad(AB)C=A(BC)\quad\ast
\]
$ABC$ 的第 $(i,j)$ 元
% 原稿在 A、B 下分别注 n、p（示 A 为 m×n、B 为 n×p）
\[
\sum_{k=1}^{n}\sum_{r=1}^{p}a_{ik}b_{kr}c_{rj}.\qedhere
\]
\end{proof}

$\forall\,A_1,A_2,\cdots,A_m$，乘积写为 $A_1A_2\cdots A_m$。

\begin{definition}[方阵的乘方]% 原稿标为「定义 7」
设 $A$ 为 $n$ 阶方阵，$A\cdot A=A^2$；
\[
\forall\,k\ge 1,\qquad A^{k}\overset{\text{def}}{=}\underbrace{A\cdot A\cdots A}_{k\ \text{个}},
\]
称为 $A$ 的\textbf{乘方}。
\end{definition}

% ===== 原稿 PDF 第 31 页 =====
% 乘方的性质、注 (1)(2)(3)、例、转置（定义 8）与转置性质及其证明

\textbf{乘方的性质}
\begin{enumerate}[label=(\arabic*)]
\item $A^{r}\cdot A^{s}=A^{r+s}$；
\item $(A^{r})^{s}=A^{rs}$\qquad$(r,s\in\Z^{+})$。
\end{enumerate}

\begin{remark}% 原稿仅标「注」，其下按 (1)(2)(3) 编号
(1) 设 $A,B$ 为 $n$ 阶方阵，一般地
\[
(A\cdot B)^{r}\ne A^{r}\cdot B^{r},
\]
\[
(A\cdot B)^{r}=\underbrace{(AB)(AB)\cdots(AB)}_{r\ \text{个}}.
\]
% 原稿第二式另起一行，仅写「=(AB)(AB)…(AB)」（下方括注 r 个），承接 (A·B)^r

(2) 若 $AB=BA$，则 $(AB)^{r}=A^{r}\cdot B^{r}$，进一步
\[
(A+B)^{m}=A^{m}+C_m^1A^{m-1}B+\cdots+C_m^{m-1}AB^{m-1}+B^{m}.
\]
\[
A\cdot I_n=I_n\cdot A=A,
\]
纯量阵 $cI_n$：
\[
A(cI_n)=(cI_n)A=cA,
\]
\[
(cI_n+A)^{m}=c^{m}I_n+C_m^1c^{m-1}A+\cdots+C_m^{m-1}cA^{m-1}+A^{m}.
\]

(3) 矩阵乘法消去律一般不成立。
\[
AB=AC\ \text{且}\ A\ne 0\ \not\Rightarrow\ B=C,
\]
% 原稿此行左侧另注（与数的情形对比）：ab=ac 且 a≠0 ⇒ b=c
「整性」一般不成立：
% 原稿「整性」二字被圈出，并以箭头连至上一行的「消去律」
\[
A\ne 0\ \text{且}\ B\ne 0\ \not\Rightarrow\ AB\ne 0.
\]
% 原稿此式下方另注（数）：a≠0, b≠0 ⇒ ab≠0
\end{remark}

\begin{example}
\[
A=B=\begin{pmatrix}0 & 1\\0 & 0\end{pmatrix},\qquad
C=\begin{pmatrix}0 & 0\\0 & 0\end{pmatrix}.
\]
\end{example}

\begin{definition}[转置矩阵、对称阵与反对称阵]% 原稿标为「定义 8」
设 $A=(a_{ij})_{m\times n}$，$A$ 的转置（$A^{T}$）$A'$ 是 $n\times m$ 阶矩阵，
\[
A'=(b_{ij})_{n\times m},
\]
其中 $b_{ij}=a_{ji}$。即 $A'$ 的第 $i$ 行就是 $A$ 的第 $i$ 列，$A'$ 的第 $j$ 列是 $A$ 的第 $j$ 行。

若 $A'=A$，则称 $A$ 为\textbf{对称阵}，
\[
A=\begin{pmatrix}1 & 3\\3 & 2\end{pmatrix};
\]
若 $A'=-A$，则称 $A$ 为\textbf{反对称阵}，
\[
A=\begin{pmatrix}0 & 1\\-1 & 0\end{pmatrix}.
\]
\end{definition}

\textbf{转置性质}
\begin{enumerate}[label=(\arabic*)]
\item $A''=A$；
\item $(A+B)'=A'+B'$；
\item $(c\cdot A)'=c\cdot A'$；
\item $(A\cdot B)'=B'\cdot A'$。
\end{enumerate}

\begin{proof}
(4) $(AB)'$ 的第 $(i,j)$ 元 $=$ $AB$ 的第 $(j,i)$ 元
\[
=A\ \text{的第}\ j\ \text{行}\cdot B\ \text{的第}\ i\ \text{列}.
\]
$B'A'$ 的第 $(i,j)$ 元 $=$ $B'$ 的第 $i$ 行 $\cdot$ $A'$ 的第 $j$ 列 $=$ $B$ 的第 $i$ 列 $\cdot$ $A$ 的第 $j$ 行。

于是 $(AB)'=B'A'$。
\end{proof}
